\begin{abstract}

    \begin{center}
        \large{% Единое место для указания названия работы
Восстановление парных зависимостей из групп непарных данных
} \\
        \large\textit{Казачков Даниил Иванович} \\[1 cm]
        \textbf{Аннотация} \\[1 cm]
    \end{center}

    Работа посвящена восстановлению зависимости между двумя наборами величин, которые
    наблюдаются группами, но исходное соответствие между их элементами утрачено. Цель работы
    состоит в разработке метода, позволяющего по таким данным одновременно оценивать скрытое
    сопоставление и общую статистическую зависимость. Сформулирована
    вероятностная модель задачи и предложены два подхода с разными параметризациями.
    Нейросетевой вариационный автоэнкодер обучается стохастическим градиентным методом, а
    модель глобальной смеси сдвигов --- модифицированным EM-алгоритмом с
    Sinkhorn-аппроксимацией E-шага и аналитическими обновлениями параметров. Исследовано
    влияние уровня шума, размера групп и объёма данных. Эксперименты на синтетических
    периодических данных показали, что оба подхода способны восстанавливать скрытые пары и
    соответствующее распределение данных. Нейросетевой подход обрабатывает новые группы без
    повторной оптимизации и показывает качество, сопоставимое с существующим методом.
    Подход с явной смесью достигает более высокой точности, когда его предположения совпадают
    со структурой данных. Установлено, что качество снижается при увеличении размера групп и
    сильном перекрытии возможных соответствий.

\end{abstract}
